\section{The Appointed Texts}
\label{sec:appointed-texts}

The Missal's texts for this Sunday are those of \latin{Dominica XXIX ``per
annum''} in the \work{Missale Romanum}, editio typica tertia (2002), which the
\work{Roman Missal}, Third Edition, renders for the dioceses of the United
States; the readings, psalm and Alleluia verse are those of no.~145 of the
\work{Lectionary for Mass}, which follows the \work{Ordo lectionum Missae} of
1981. The Latin of the Missal, the approved English of the Missal and the
approved English of the Lectionary are all under copyright, and none of them
is reproduced here. Each biblical element is given at its canonical verses in
the Douay--Rheims version as revised by Bishop Challoner, which translates the
Clementine Vulgate and numbers the psalms as the Vulgate does. That English is
a study text and is not the wording proclaimed at Mass. Words printed in
brackets and italics belong to the verse but not to the appointed text. Each
oration is identified by the opening words of its Latin and described, and no
English is supplied for it. The three antiphons and the Alleluia verse adapt
Scripture: the verses they draw on are printed beside them as their sources and
not as their words, and no English is composed for any of them.

\Needspace{8\baselineskip}
\subsection{Entrance Antiphon}
\label{proper-entrance-antiphon}

\properlocus{\latin{Ego clamavi, quoniam exaudisti me, Deus}. The Missal cites
Psalm 16:6 and 8 in the Vulgate numbering, which is Psalm 17:6 and 8 in the
Hebrew, and prints \latin{Cf.}}

\noindent A prayer in two movements. The speaker has called on God because God
hears him, and asks him to bend his ear and hear his words; then he asks to be
guarded like the pupil of the eye and protected under the shadow of God's
wings. The antiphon takes v.~6 without the Clementine Vulgate's \latin{mihi},
passes over v.~7, and gives v.~8 without its opening phrase, \latin{A
resistentibus dexterae tuae}, adding the vocative \latin{Domine} after
\latin{Custodi me}. It follows the Clementine wording of v.~6, \latin{Ego
clamavi, quoniam exaudisti me}, rather than the \work{Nova Vulgata}'s, and its
added vocative stands in neither Latin psalter. No English is printed for the
antiphon. The verses it draws on follow.

\begin{englishwitness}{Douay--Rheims (Challoner), Ps 16:6--8, the antiphon's source}
\textsuperscript{6}I have cried to thee, for thou, O God, hast heard me: O
incline thy ear unto me, and hear my words. \notread{\textsuperscript{7}Shew
forth thy wonderful mercies; thou who savest them that trust in thee.}
\textsuperscript{8}\notread{From them that resist thy right hand} keep me, as
the apple of thy eye. Protect me under the shadow of thy wings.
\end{englishwitness}

\Needspace{8\baselineskip}
\subsection{Collect}
\label{proper-collect}

\properlocus{\latin{Omnipotens sempiterne Deus, fac nos tibi semper}. Long
conclusion.}

\noindent Addressed to the almighty and eternal God, a petition with two
members: that he make the assembly always (\latin{semper}) hold its will
devoted to him, and serve his majesty (\latin{maiestati tuae \ldots\ servire})
with an undivided heart (\latin{sincero corde}). No English is printed.

\Needspace{8\baselineskip}
\subsection{First Reading}
\label{proper-first-reading}

\properlocus{Isaiah 45:1, 4--6. Lectionary no.~145. The reading is in two
parts: v.~1, then vv.~4--6; vv.~2--3 are left out between them. The
\work{Ordo}'s Latin begins \latin{Haec dicit Dominus de uncto meo Cyro}.}

\begin{englishwitness}{Douay--Rheims (Challoner), Is 45:1--6}
\textsuperscript{1}Thus saith the Lord to my anointed Cyrus, whose right hand
I have taken hold of, to subdue nations before his face, and to turn the backs
of kings, and to open the doors before him, and the gates shall not be shut.
\notread{\textsuperscript{2}I will go before thee, and will humble the great
ones of the earth: I will break in pieces the gates of brass, and will burst
the bars of iron. \textsuperscript{3}And I will give thee hidden treasures, and
the concealed riches of secret places: that thou mayest know that I am the Lord
who call thee by thy name, the God of Israel.}
\textsuperscript{4}For the sake of my servant Jacob, and Israel my elect, I
have even called thee by thy name: I have made a likeness of thee, and thou
hast not known me. \textsuperscript{5}I am the Lord, and there is none else:
there is no God besides me: I girded thee, and thou hast not known me:
\textsuperscript{6}That they may know who are from the rising of the sun, and
they who are from the west, that there is none besides me. I am the Lord, and
there is none else:
\end{englishwitness}

\noindent Two words of this study text differ from the Lectionary's Latin.
In v.~1 the Douay--Rheims ``my anointed'' translates the Clementine's
\latin{christo meo}; the \work{Ordo} reads \latin{uncto meo}, the same title
in a Latin word that does not transliterate the Greek. In v.~4 ``I have made a
likeness of thee'' translates the Clementine's \latin{assimilavi te}; the
Lectionary's English follows the sense of the Hebrew, which the King James
Version renders ``I have surnamed thee'', and does not speak of a likeness.
Both points matter below.

\Needspace{8\baselineskip}
\subsection{Responsorial Psalm}
\label{proper-responsorial-psalm}

\properlocus{Psalm 96 (95):1 and 3, 4--5, 7--8, 9--10a and 10c; the response is
taken from v.~7b. Lectionary no.~145. Verses~2 and~6 and the middle clause of
v.~10 are not sung, and the title is not part of the first strophe.}

\noindent The \work{Ordo} prints the response in the Latin \latin{Afferte
Domino gloriam et potentiam}, glory and might, which is the \work{Nova
Vulgata}'s wording of v.~7b; the Clementine and the Greek have glory and
honour. The response on the bishops' conference's page of readings for the day
names glory and honour, while its own strophe renders the same half-verse with
a different second noun. Nothing below rests on the second noun of the
response.

\begin{englishwitness}{Douay--Rheims (Challoner), Ps 95:1--10}
\textsuperscript{1}\notread{A canticle for David himself, when the house was
built after the captivity.} Sing ye to the Lord a new canticle: sing to the
Lord, all the earth. \notread{\textsuperscript{2}Sing ye to the Lord and bless
his name: shew forth his salvation from day to day.}
\textsuperscript{3}Declare his glory among the Gentiles: his wonders among all
people.

\textsuperscript{4}For the Lord is great, and exceedingly to be praised: he is
to be feared above all gods. \textsuperscript{5}For all the gods of the
Gentiles are devils: but the Lord made the heavens. \notread{\textsuperscript{6}Praise
and beauty are before him: holiness and majesty in his sanctuary.}

\textsuperscript{7}Bring ye to the Lord, O ye kindreds of the Gentiles, bring
ye to the Lord glory and honour: \textsuperscript{8}Bring to the Lord glory
unto his name. Bring up sacrifices, and come into his courts:

\textsuperscript{9}Adore ye the Lord in his holy court. Let all the earth be
moved at his presence. \textsuperscript{10}Say ye among the Gentiles, the Lord
hath reigned. \notread{For he hath corrected the world, which shall not be
moved:} he will judge the people with justice.
\end{englishwitness}

\noindent In v.~5 the Douay--Rheims ``devils'' follows the Clementine
\latin{daemonia} and the Greek; the Hebrew-based versions read ``idols'', and
the \work{Nova Vulgata}, as the Holy See publishes it, has \latin{inania},
empty things. The Greek and Latin
title, which the Hebrew psalm does not carry, ties the psalm to the house
rebuilt after the captivity.

\Needspace{8\baselineskip}
\subsection{Second Reading}
\label{proper-second-reading}

\properlocus{1~Thessalonians 1:1--5b. Lectionary no.~145. The reading ends
within v.~5; its last clause is not read.}

\begin{englishwitness}{Douay--Rheims (Challoner), 1~Thes 1:1--5}
\textsuperscript{1}Paul and Sylvanus and Timothy to the church of the
Thessalonians: in God the Father and in the Lord Jesus Christ.
\textsuperscript{2}Grace be to you and peace. We give thanks to God always for
you all: making a remembrance of you in our prayers without ceasing,
\textsuperscript{3}Being mindful of the work of your faith and labour and
charity: and of the enduring of the hope of our Lord Jesus Christ before God
and our Father. \textsuperscript{4}Knowing, brethren, beloved of God, your
election: \textsuperscript{5}For our gospel hath not been unto you in word
only, but in power also: and in the Holy Ghost and in much fulness\notread{, as
you know what manner of men we have been among you for your sakes.}
\end{englishwitness}

\noindent The Douay--Rheims sets the greeting ``Grace be to you and peace'' at
the beginning of v.~2; the division the Lectionary follows keeps it at the end
of v.~1. The reading includes it either way.

\Needspace{8\baselineskip}
\subsection{Alleluia Verse}
\label{proper-acclamation}

\properlocus{Philippians 2:15d, 16a. Lectionary no.~145; printed without
\latin{Cf.} The \work{Ordo}'s Latin verse is \latin{Lucetis sicut luminaria in
mundo, verbum vitae continentes}.}

\noindent The verse joins the last colon of v.~15 to the first of v.~16. The
\work{Ordo} gives both in the indicative, as the Vulgate does, leaving out the
relative \latin{inter quos}, ``among whom''. The United States verse on the
bishops' conference's page phrases them as an exhortation to shine and to hold
on to the word of life. The Greek verb can be read either way, and its
participle can mean holding fast or holding out; nothing below rests on the
mood of the English verb. No English is printed for the verse; its source
follows, with the words outside the verse in brackets.

\begin{englishwitness}{Douay--Rheims (Challoner), Phil 2:15--16, the verse's source}
\notread{\textsuperscript{15}That you may be blameless and sincere children of
God, without reproof, in the midst of a crooked and perverse generation: among
whom} you shine as lights in the world. \textsuperscript{16}Holding forth the
word of life \notread{to my glory in the day of Christ: because I have not run
in vain, nor laboured in vain.}
\end{englishwitness}

\Needspace{8\baselineskip}
\subsection{Gospel}
\label{proper-gospel}

\properlocus{Matthew 22:15--21. Lectionary no.~145. The \work{Ordo} opens the
reading \latin{In illo tempore} and names Jesus in v.~15.}

\begin{englishwitness}{Douay--Rheims (Challoner), Mt 22:15--21}
\textsuperscript{15}Then the Pharisees going, consulted among themselves how to
insnare him in his speech. \textsuperscript{16}And they sent to him their
disciples with the Herodians, saying: Master, we know that thou art a true
speaker and teachest the way of God in truth. Neither carest thou for any man:
for thou dost not regard the person of men. \textsuperscript{17}Tell us
therefore what dost thou think? Is it lawful to give tribute to Caesar, or
not? \textsuperscript{18}But Jesus knowing their wickedness, said: Why do you
tempt me, ye hypocrites? \textsuperscript{19}Shew me the coin of the tribute.
And they offered him a penny. \textsuperscript{20}And Jesus saith to them:
Whose image and inscription is this? \textsuperscript{21}They say to him:
Caesar's. Then he saith to them: Render therefore to Caesar the things that are
Caesar's; and to God, the things that are God's.
\end{englishwitness}

\noindent The ``penny'' is a denarius. ``Tribute'' renders the Latin
\latin{censum} and the Greek \textit{k\=ensos}, the census, which Aquinas
explains as the tribute paid for each head (\latin{Census erat tributum quod
pro capite dabatur}); the coin is ``the coin of the tribute''.

\Needspace{8\baselineskip}
\subsection{Prayer over the Offerings}
\label{proper-prayer-over-offerings}

\properlocus{\latin{Tribue nos, Domine, quaesumus, donis tuis}. Short
conclusion.}

\noindent A petition about the manner of service at the altar: that those who
offer may serve the Lord's own gifts with a free mind (\latin{libera mente
servire}), so that his grace, purifying them, may make them clean through the
same mysteries in which they minister (\latin{iisdem quibus famulamur
mysteriis emundemur}). No English is printed.

\Needspace{8\baselineskip}
\subsection{Communion Antiphon, first option}
\label{proper-communion-antiphon-a}

\properlocus{\latin{Ecce oculi Domini super timentes eum}. The Missal cites
Psalm 32:18--19 in the Vulgate numbering, which is Psalm 33:18--19 in the
Hebrew, and prints \latin{Cf.}}

\noindent The Lord's eyes rest on those who fear him and hope in his mercy, to
rescue their souls from death and to feed them in famine. The antiphon reads
\latin{timentes} where both the Clementine and the \work{Nova Vulgata} read
\latin{metuentes}; otherwise it follows the two verses. The verses it draws on
follow, as its source.

\begin{englishwitness}{Douay--Rheims (Challoner), Ps 32:18--19, the antiphon's source}
\textsuperscript{18}Behold the eyes of the Lord are on them that fear him: and
on them that hope in his mercy. \textsuperscript{19}To deliver their souls from
death; and feed them in famine.
\end{englishwitness}

\Needspace{8\baselineskip}
\subsection{Communion Antiphon, second option}
\label{proper-communion-antiphon-b}

\properlocus{\latin{Filius hominis venit}. The Missal cites Mark 10:45 and
prints no \latin{Cf.}}

\noindent The Son of Man came to give his life as a ransom for many. The
antiphon keeps only the end of the verse: it omits the opening \latin{Nam et},
``for even'', and the whole contrast between being served and serving, so that
the Son of Man simply came to give his life. No English is printed for the
antiphon; its source follows, with the words it omits in brackets.

\begin{englishwitness}{Douay--Rheims (Challoner), Mk 10:45, the antiphon's source}
\notread{For} the Son of man \notread{also is not} come \notread{to be
ministered unto: but to minister and} to give his life a redemption for many.
\end{englishwitness}

\Needspace{8\baselineskip}
\subsection{Prayer after Communion}
\label{proper-prayer-after-communion}

\properlocus{\latin{Fac nos, quaesumus, Domine, caelestium rerum}. Short
conclusion.}

\noindent A petition that frequent sharing in heavenly things
(\latin{caelestium rerum frequentatione}) may make the assembly advance, so
that it is aided by temporal benefits and instructed (\latin{erudiamur}) by
eternal ones. No English is printed.
